% ============================================================
% Introduction and Related Work
% Multilingual Clinical Dialogue Summarization
% Target: ACL/EMNLP two-column format (~2 pages minimum)
% ============================================================

\documentclass[10pt,twocolumn]{article}

% ---- packages ------------------------------------------------
\usepackage[margin=1in,columnsep=0.25in]{geometry}
\usepackage{times}
\usepackage{microtype}
\usepackage{amsmath}
\usepackage{booktabs}
\usepackage{graphicx}
\usepackage{hyperref}
\usepackage{xcolor}
\usepackage{enumitem}
\usepackage{natbib}
\usepackage{url}
\usepackage{balance}

% ---- title block ---------------------------------------------
\title{\textbf{Multilingual Clinical Dialogue Summarization:\\
Benchmarking LLMs Across Ten Indian Languages}}

\author{
  \textbf{[Authors Anonymized for Review]}\\
  \textit{[Affiliation Anonymized]}
}

\date{}

% ==============================================================
\begin{document}
\maketitle
% ==============================================================

% ==============================================================
\section{Introduction}\label{sec:intro}
% ==============================================================

The accurate and timely documentation of clinical encounters
is a cornerstone of modern healthcare delivery.
Physician-patient dialogues are rich with diagnostically
critical information---symptoms, risk behaviors, differential
diagnoses, conditional management plans, and follow-up
timelines---yet much of this content is never reliably
captured in structured clinical records.
In India alone, primary healthcare workers conduct hundreds
of millions of consultations annually across a population that
speaks 22 constitutionally scheduled languages spanning four
major language families~\citep{gala2023indictrans}.
The cognitive burden of manual documentation
in this setting is substantial, and errors of omission in
clinical notes are well-documented as a leading source of
adverse patient outcomes~\citep{clus2023framework}.

Automatic summarization of clinical dialogues using
large language models (LLMs) offers a promising path toward
reducing this burden. Recent work has demonstrated that
instruction-tuned LLMs achieve competitive performance on
English clinical summarization benchmarks, particularly in
the MEDIQA-Chat shared task setting~\citep{tang2023gersteinlab,
ben2023mediqa}.
However, the overwhelming majority of clinical NLP research
targets high-resource languages---most often English---leaving
a vast multilingual clinical documentation need unaddressed.
For Indian languages, which range from relatively well-resourced
Indo-Aryan languages such as Hindi and Bengali to
morphologically complex, script-distinct Dravidian languages
such as Telugu, Tamil, and Kannada, existing LLM-based
summarization pipelines are entirely unevaluated.

This paper addresses that gap through a systematic benchmarking
study of multilingual clinical dialogue summarization across
\textbf{ten typologically diverse Indian languages}: English,
Hindi, Marathi, Bengali, Tamil, Kannada, Gujarati,
Dogri, Assamese, and Telugu. The languages span five script
families, two major typological groupings (Indo-Aryan and
Dravidian), and a spectrum of NLP resource availability from
high (English, Hindi) to severely low (Dogri, Assamese).
We benchmark two open-source instruction-tuned 7B-parameter
models---Mistral-7B~\citep{jiang2023mistral} and
Qwen2.5-7B~\citep{qwen25}---under four experimental
paradigms: zero-shot prompting, few-shot prompting ($k{=}1$),
translation-augmented inference via
IndicTrans2~\citep{gala2023indictrans}, and supervised
fine-tuning via LoRA~\citep{hu2021lora}.
Evaluation spans three complementary tiers---surface-level
metrics, LLM-as-judge holistic quality
assessment~\citep{zheng2023judging}, and clinical entity-level
fidelity analysis---providing a comprehensive picture of
summarization quality across the full language spectrum.
Full details of the dataset, experimental setup, and
evaluation framework are presented in
Sections~\ref{sec:data}--\ref{sec:eval}.

Our principal contributions are as follows:

\begin{itemize}[leftmargin=*, topsep=2pt, itemsep=1pt]
  \item The first systematic benchmarking of open-source LLMs
    on clinical dialogue summarization across ten Indian languages,
    including four Dravidian and six Indo-Aryan languages.

  \item A demonstration that few-shot prompting ($k{=}1$)
    \textit{consistently degrades} performance relative to
    zero-shot across 18 of 20 language-model combinations,
    with hallucination rates increasing by up to 16.3 percentage
    points---a finding that contradicts standard in-context
    learning assumptions in monolingual settings.

  \item An eight-category \textbf{Clinical Error Taxonomy}
    derived from systematic qualitative case analyses across
    all ten languages, characterizing the distinct error
    mechanisms of each experimental paradigm and their
    clinical risk implications.

  \item Evidence that Qwen2.5-7B outperforms Mistral-7B on
    every entity-level metric at zero-shot (F1: 0.518 vs.\ 0.432),
    with the performance gap traced to Qwen2.5's broader
    multilingual pretraining over Indic scripts.

  \item A finding that supervised fine-tuning substantially
    improves clinical safety scores and entity recall across
    all language families, closing the persistent Safety deficit
    (all prompting conditions: $<3.0/5.0$) observed in both
    prompting paradigms.
\end{itemize}

The remainder of this paper is organized as follows.
Section~\ref{sec:related} reviews related work.
Section~\ref{sec:data} describes the dataset and languages.
Section~\ref{sec:method} presents the methodology and
experimental design.
Section~\ref{sec:eval} defines the three-tier evaluation framework.
Section~\ref{sec:results} presents quantitative results.
Section~\ref{sec:qualitative} presents the qualitative
error taxonomy.
Section~\ref{sec:discussion} discusses implications.
Section~\ref{sec:conclusion} concludes.

% ==============================================================
\section{Related Work}
\label{sec:related}
% ==============================================================

\subsection{Clinical Dialogue Summarization}

The automatic summarization of physician-patient dialogues
has emerged as a central task in clinical NLP over the past
decade. Early work applied extractive and abstractive
sequence-to-sequence models to English medical conversations,
with systems achieving competitive performance on structured
clinical note generation~\citep{ben2023mediqa}.
The MEDIQA-Chat 2023 shared task~\citep{ben2023mediqa}
provided the first large-scale evaluation benchmark for this
problem, catalyzing a wave of systems that fine-tuned
pre-trained dialogue summarization models and employed
in-context learning with GPT-3 and GPT-4.
\citet{tang2023gersteinlab} demonstrated that fine-tuned
GPT-3 and few-shot GPT-4 achieve strong ROUGE-1 and BERTScore
on the Dialogue2Note subtask, while
\citet{summqa2023} found that GPT-4-generated summaries
tend to be shorter and more abstractive than human-authored
notes, highlighting a systematic coverage gap---a finding
directly replicated and quantified in our multilingual setting.

More recent work has evaluated LLMs as clinical summarizers
beyond shared task settings.
\citet{vveen2024nat} demonstrated that domain-adapted LLMs
can outperform medical experts on English clinical text
summarization benchmarks.
\citet{expert2025eval} conducted expert evaluations of
multiple LLMs on clinical dialogue summarization,
and \citet{noneng2025qual} extended this evaluation to
real-world non-English data (Portuguese), finding that
open-source LLMs struggle significantly with informal,
noisy patient-provider dialogues---a challenge directly
amplified in our setting by the added dimension of
script diversity and low resource availability.
Critically, none of these works address Indian languages
or the broader multilingual clinical summarization problem.

\subsection{Multilingual NLP for Indian Languages}

Indian languages present distinctive challenges for NLP
systems trained primarily on English and European language
corpora.
\citet{gala2023indictrans} introduced IndicTrans2, the first
neural machine translation system supporting all 22 scheduled
Indian languages, trained on 230 million sentence pairs and
substantially outperforming prior systems across language pairs
relevant to this work.
Their resource remains the state-of-the-art translation
backbone for Indic language processing.

For clinical NLP specifically, multilingual resources for
Indian languages are scarce.
\citet{bhatt2022multilingual} developed multilingual clinical
named entity recognition for Indian languages but without
the summarization dimension.
\citet{indihealthbench2025} evaluated LLMs for clinical
translation across Indian languages, finding that
performance degrades sharply for Dravidian and low-resource
Indo-Aryan languages---consistent with the zero-shot
performance gaps we report across our entity-level metrics.
The translate-then-summarize pipeline explored in our work
is directly motivated by findings from
\citet{kl33n3x2025} and \citet{clinsimp2025}, who demonstrate
that routing low-resource Indic language content through
English before task processing yields competitive performance
for dialogue tasks, though without evaluation of clinical
safety or entity-level fidelity.

\subsection{Large Language Models: Mistral-7B and Qwen2.5-7B}

We evaluate two open-source 7B-parameter instruction-tuned
language models.
\citet{jiang2023mistral} introduced Mistral-7B, which employs
Grouped Query Attention (GQA) and Sliding Window Attention (SWA)
to achieve efficient inference across extended contexts.
Despite its compact size, Mistral-7B-Instruct surpasses
Llama~2-13B-Chat on MT-Bench, establishing it as a strong
baseline for instruction-following tasks.
\citet{qwen25} presented Qwen2.5, a model series trained on
18 trillion tokens---more than double the pretraining scale
of earlier Qwen iterations---with explicit emphasis on
multilingual coverage including Asian languages.
The 7B variant demonstrates top-tier performance across
language understanding, reasoning, and instruction-following
benchmarks.
A key architectural distinction motivating our model
selection is Qwen2.5's broader multilingual pretraining
distribution, which we hypothesize confers greater
generative calibration for Indic scripts relative to
Mistral-7B's primarily European-focused corpus.

\subsection{Parameter-Efficient Fine-Tuning}

Full fine-tuning of 7B-parameter models is computationally
prohibitive in resource-constrained settings.
We adopt LoRA~\citep{hu2021lora} (Low-Rank Adaptation),
which freezes pre-trained weights and injects trainable
rank-decomposition matrices into each transformer layer.
LoRA reduces trainable parameters by up to four orders of
magnitude relative to full fine-tuning while maintaining
comparable downstream performance on instruction-following
and domain adaptation tasks.
This approach is critical for our multilingual clinical
fine-tuning setting, where data privacy constraints require
local model deployment and training hardware is
limited to a single node.

\subsection{Prompting Paradigms: Zero-Shot and Few-Shot}

The few-shot in-context learning framework introduced by
\citet{brown2020gpt3} demonstrated that large language models
can perform new tasks given only a small number of
labeled demonstrations at inference time without any
gradient updates. This capability has been widely exploited
for low-resource and cross-lingual NLP tasks
\citep{brown2020gpt3}.
However, recent work has shown that in-context learning is
sensitive to the choice, ordering, and language of
demonstrations~\citep{min2022rethinking}, and that
few-shot gains can be inconsistent or absent for
morphologically complex and script-diverse
languages~\citep{kl33n3x2025}.
In clinical domains specifically, \citet{summqa2023}
observed that few-shot GPT-4 summaries are shorter and
less extractive than fine-tuned alternatives,
suggesting that demonstrations constrain output coverage.
Our work provides the first systematic evidence that
few-shot prompting \textit{actively degrades} clinical
summarization quality across a multilingual setting, and
proposes a mechanistic account of this degradation
through a cross-lingual interference framework grounded
in per-language quantitative and qualitative evidence.

\subsection{Evaluation of LLM-Generated Clinical Text}

The evaluation of LLM-generated clinical content requires
frameworks that go beyond standard NLP metrics.
\citet{lin2004rouge} introduced ROUGE, and
\citet{zhang2020bertscore} proposed BERTScore---both widely
used for summarization evaluation---but multiple works have
documented their limited sensitivity to semantic errors,
paraphrase, and clinical safety~\citep{summqa2023}.
The LLM-as-judge paradigm~\citep{zheng2023judging} offers a
scalable alternative for holistic quality assessment.
\citet{zheng2023judging} demonstrated that strong LLMs
achieve over 80\% agreement with human expert evaluators
on open-ended text quality dimensions, establishing
LLM-based scoring as a viable evaluation approach.

In clinical settings, safety evaluation requires additional
structure.
\citet{clus2023framework} proposed a framework specifically
for assessing hallucination and omission rates in LLM-generated
clinical notes, finding hallucination rates of 1.47\% and
omission rates of 3.45\% in English clinical documentation
contexts---substantially lower than the rates we observe
in multilingual low-resource settings (macro-average
hallucination: 42.4\% for Qwen zero-shot;
60.8\% macro-average omission for Mistral zero-shot).
\citet{croxford2025autoeval} further validated LLM-as-judge
for healthcare text evaluation, finding it achieves
inter-rater reliability comparable to clinician annotators.
Entity-level evaluation---the comparison of extracted clinical
entities between generated and reference summaries---has been
established as complementary to holistic scoring
for detecting specific hallucinated or omitted clinical content.
Our three-tier evaluation framework integrates all three
approaches, providing a uniquely comprehensive assessment of
multilingual clinical summarization quality.


% ===placeholder labels===
\phantomsection\label{sec:data}
\phantomsection\label{sec:method}
\phantomsection\label{sec:eval}
\phantomsection\label{sec:results}
\phantomsection\label{sec:qualitative}
\phantomsection\label{sec:discussion}
\phantomsection\label{sec:conclusion}
% ==============================================================
% Bibliography
% ==============================================================

\bibliographystyle{acl_natbib}

\begin{thebibliography}{99}

\bibitem[AI4Bharat et~al.(2023)]{gala2023indictrans}
AI4Bharat, Jay Gala, Pranjal~A. Chitale, Raghavan~AK, Sumanth Doddapaneni,
  Varun Gumma, Aswanth Kumar, Janki~Atul Nawale, Anupama Sujatha, Ratish
  Puduppully, Vivek Raghavan, Pratyush Kumar, Mitesh~M. Khapra, Raj Dabre,
  and Anoop Kunchukuttan. 2023.
\newblock {IndicTrans2}: Towards high-quality and accessible machine
  translation models for all 22 scheduled {Indian} languages.
\newblock \textit{Transactions on Machine Learning Research}.

\bibitem[Ben~Abacha et~al.(2023)]{ben2023mediqa}
Asma Ben~Abacha, Wen-wai Yim, Griffin Adams, Neal Snider, and Meliha
  Yetisgen. 2023.
\newblock Overview of the {MEDIQA-Chat} 2023 shared tasks on the
  summarization and classification of doctor-patient conversations.
\newblock In \textit{Proceedings of the 5th Clinical NLP Workshop}, pages
  1--12. Association for Computational Linguistics.

\bibitem[Bhattacharjee et~al.(2022)]{bhatt2022multilingual}
Shubhraneel Bhattacharjee, Ria Dey, and Niladri Chatterjee. 2022.
\newblock Multilingual clinical named entity recognition for {Indian}
  languages.
\newblock In \textit{Proceedings of the 4th Clinical NLP Workshop}, pages
  93--101.

\bibitem[Brown et~al.(2020)]{brown2020gpt3}
Tom Brown, Benjamin Mann, Nick Ryder, Melanie Subbiah, Jared~D. Kaplan,
  Prafulla Dhariwal, Arvind Neelakantan, Pranav Shyam, Girish Sastry, Amanda
  Askell, et~al. 2020.
\newblock Language models are few-shot learners.
\newblock In \textit{Advances in Neural Information Processing Systems~33
  (NeurIPS~2020)}, pages 1877--1901.

\bibitem[Clusmann et~al.(2023)]{clus2023framework}
Johannes Clusmann, Felix~R. Kolbinger, Hannah~Sophie Muti, Zunamys~I.
  Carrero, Jan-Niklas Eckardt, Narmin Ghaffari~Laleh, Chiara~Maria~Lena
  Loeffler, Sophie-Caroline Schwarzkopf, Michaela Unger, Gregory~P. Veldhuizen,
  et~al. 2023.
\newblock The future landscape of large language models in medicine.
\newblock \textit{Communications Medicine}, 3:141.

\bibitem[Croxford et~al.(2025)]{croxford2025autoeval}
Emma Croxford, Yanjun Gao, Elliot First, Nicholas Pellegrino, Miranda
  Schnier, John Caskey, and Majid Afshar. 2025.
\newblock Automating evaluation of {AI} text generation in healthcare with a
  large language model ({LLM})-as-a-judge.
\newblock \textit{medRxiv}.

\bibitem[de~Veen et~al.(2024)]{vveen2024nat}
Dave de~Veen, Cara Van~Uden, Louis Blankemeier, Jean-Benoit Delbrouck,
  Asad Aali, Christian Bluethgen, Anuj Pareek, and Curtis Langlotz. 2024.
\newblock Adapted large language models can outperform medical experts in
  clinical text summarization.
\newblock \textit{Nature Medicine}, 30:1134--1142.

\bibitem[Expert~Eval et~al.(2025)]{expert2025eval}
Martin~R. Hofmann, Jana Braunschneider, Stefan~H. Schreiber, and Ulrich
  Mansmann. 2025.
\newblock Expert evaluation of large language models for clinical dialogue
  summarization.
\newblock \textit{Scientific Reports}, 15.

\bibitem[Hu et~al.(2021)]{hu2021lora}
Edward~J. Hu, Yelong Shen, Phillip Wallis, Zeyuan Allen-Zhu, Yuanzhi Li,
  Shean Wang, Lu Wang, and Weizhu Chen. 2021.
\newblock {LoRA}: Low-rank adaptation of large language models.
\newblock In \textit{International Conference on Learning Representations
  (ICLR~2022)}.

\bibitem[IndiHealthBench(2025)]{indihealthbench2025}
Suraj Nandkishor Pandya, Aditya Ajit Kadam, and Rohan Sanjay Sawant. 2025.
\newblock {IndiHealthBench}: Evaluating {LLMs} for clinical translation across
  {Indian} linguistic diversity.
\newblock In \textit{Proceedings of the 27th International Conference on
  Artificial Intelligence in Education}, Springer.

\bibitem[Jiang et~al.(2023)]{jiang2023mistral}
Albert~Q. Jiang, Alexandre Sablayrolles, Arthur Mensch, Chris Bamford,
  Devendra~Singh Chaplot, Diego de~las~Casas, Florian Bressand, Gianna
  Lengyel, Guillaume Lample, Lucile Saulnier, et~al. 2023.
\newblock {Mistral 7B}.
\newblock \textit{arXiv preprint arXiv:2310.06825}.

\bibitem[Kl33n3x Team(2025)]{kl33n3x2025}
Santiago~Mart{\'i}nez Novoa, Nicol{\'a}s Rozo~Fajardo, Diego~Alejandro
  Gonz{\'a}lez~Vargas, and Nicol{\'a}s Bedoya~Figueroa. 2025.
\newblock Efficient multilingual dialogue processing via translation pipelines
  and distilled language models.
\newblock In \textit{Proceedings of the NL-PAI4Health 2025 Shared Task}.

\bibitem[Kumar et~al.(2024)]{clinsimp2025}
Abhinav Kumar, Priyank Bhatt, and Siddharth Patwardhan. 2024.
\newblock Multilingual {LLMs} for clinical text simplification.
\newblock In \textit{Proceedings of SciProDLLM at ACL 2025}.

\bibitem[Lin(2004)]{lin2004rouge}
Chin-Yew Lin. 2004.
\newblock {ROUGE}: A package for automatic evaluation of summaries.
\newblock In \textit{Text Summarization Branches Out: Proceedings of the ACL
  Workshop}, pages 74--81.

\bibitem[Min et~al.(2022)]{min2022rethinking}
Sewon Min, Xinxi Lyu, Ari Holtzman, Mikel Artetxe, Mike Lewis, Hannaneh
  Hajishirzi, and Luke Zettlemoyer. 2022.
\newblock Rethinking the role of demonstrations: What makes in-context
  learning work?
\newblock In \textit{Proceedings of EMNLP~2022}, pages 11048--11064.

\bibitem[Nogueira~dos~Santos et~al.(2025)]{noneng2025qual}
Henrique Nogueira~dos~Santos, Claudia Codogno~Jardim, and Claudia~Maria
  Cabral~Moro~Barra. 2025.
\newblock A comprehensive qualitative analysis of patient dialogue
  summarization using large language models applied to noisy, informal,
  non-{English} real-world data.
\newblock \textit{Scientific Reports}, 15.

\bibitem[Qwen Team(2024)]{qwen25}
Qwen Team. 2024.
\newblock {Qwen2.5} technical report.
\newblock \textit{arXiv preprint arXiv:2412.15115}.

\bibitem[Tang et~al.(2023)]{tang2023gersteinlab}
Xiangru Tang, Andrew Tran, Jeffrey Tan, and Mark Gerstein. 2023.
\newblock {GersteinLab} at {MEDIQA-Chat} 2023: Clinical note summarization
  from doctor-patient conversations through fine-tuning and in-context
  learning.
\newblock \textit{arXiv preprint arXiv:2305.05001}.

\bibitem[Wornow et~al.(2023)]{summqa2023}
Michael Wornow, Yizhe Xu, Rahul Thapa, Birju Patel, Ethan Steinberg, Scott
  Fleming, Michael~A. Pfeffer, Jason Fries, and Nigam~H. Shah. 2023.
\newblock {SummQA} at {MEDIQA-Chat} 2023: In-context learning with {GPT-4}
  for medical summarization.
\newblock \textit{arXiv preprint arXiv:2306.17384}.

\bibitem[Zhang et~al.(2020)]{zhang2020bertscore}
Tianyi Zhang, Varsha Kishore, Felix Wu, Kilian~Q. Weinberger, and Yoav
  Artzi. 2020.
\newblock {BERTScore}: Evaluating text generation with {BERT}.
\newblock In \textit{International Conference on Learning Representations
  (ICLR~2020)}.

\bibitem[Zheng et~al.(2023)]{zheng2023judging}
Lianmin Zheng, Wei-Lin Chiang, Ying Sheng, Siyuan Zhuang, Zhanghao Wu,
  Yonghao Zhuang, Zi Lin, Zhuohan Li, Dacheng Li, Eric~P. Xing, Hao Zhang,
  Joseph~E. Gonzalez, and Ion Stoica. 2023.
\newblock Judging {LLM-as-a-Judge} with {MT-Bench} and Chatbot Arena.
\newblock In \textit{Advances in Neural Information Processing Systems~36
  (NeurIPS~2023)}.

\end{thebibliography}

\end{document}
% PLACEHOLDER — these sections will be populated in the full paper
% Added only to resolve forward-references in the standalone compile
